\chapter{机器在夜里做什么}
% 完整版：chapters/20_sleep.tex

睡眠不是关机，而是切换模式。从电台就能看出来。在深睡眠中，几乎所有电台都调低了：乙酰胆碱、去甲肾上腺素、血清素。在做梦的睡眠中，乙酰胆碱回到白天的水平，去甲肾上腺素和血清素几乎沉默，而通往肌肉的输出被刹车关掉——所以我们睡着时不会真的跑去追梦里要逃离的东西。这些全都是测量到的。

\section*{什么时候该睡了}

想象一个电容，在你醒着的一整天里充电，夜里放电。电荷就是“睡眠压力”。电荷达到上阈值，你就睡着；降到下阈值，你就醒来。而这两个阈值本身会随着上一章的昼夜节律平缓地上下移动。这样一个简单的电路，自己就会落到清醒十六小时、睡眠八小时。它还解释了为什么熬了一夜之后，我们不是睡两倍长，而是睡得更深：充满的电容放电更快。“电荷”的候选者是一种在大脑工作时积累起来的物质；是否正是它，目前还是假说。

\pencilfig{120}{%
% figures/sleep_{S,H,L}.dat from models/sleep_models.py, t in hours
\begin{scope}[x=0.1cm, y=2.6cm]
\foreach \a/\b in {0/4.3,21.2/29.3,45.4/53.4,69.5/72}{\fill[pcBlue, opacity=0.1] (\a,0) rectangle (\b,1);}
\draw[pen] (0,0) -- (72,0);
\draw[penlite=pcRed, densely dashed] plot file {figures/sleep_H.dat};
\draw[penlite=pcGreen, densely dashed] plot file {figures/sleep_L.dat};
\draw[pcBlue, line width=1.0pt] plot file {figures/sleep_S.dat};
\foreach \t/\l in {0/0,24/一天,48/两天,72/三天}{\draw[pcGray, line width=0.5pt] (\t,0) -- (\t,-0.04); \node[lbl, anchor=base] at (\t,-0.16) {\l};}
\node[lbl, text=pcBlue] at (25.25,0.93) {睡眠}; \node[lbl, text=pcBlue] at (49.4,0.93) {睡眠};
\end{scope}
\node[lbl, text=pcRed, anchor=west] at (7.3,2.0) {入睡};
\node[lbl, text=pcGreen!70!black, anchor=west] at (7.3,0.62) {醒来};
\node[lbl, text=pcBlue, anchor=west] at (7.3,1.35) {睡眠压力};
}{睡眠压力白天积累、夜里消退；到上阈值时入睡，到下阈值时醒来，而两个阈值都随昼夜起伏。}

\section*{慢速跷跷板}

在深睡眠中，大脑的整片区域同步地时而开启、时而关闭——频率低于每秒一次。这是我们熟悉的电路：一群自我兴奋、又会疲劳的神经元——也就是振荡器。白天，乙酰胆碱抑制神经元的疲劳，同一群神经元平稳工作；夜里乙酰胆碱少了，网络就开始摇摆。

\pencilfig{20}{\pencilart{20}}{夜里网络缓慢摇摆：一会儿几乎所有神经元都在工作，一会儿几乎全都沉默。}

\section*{快进}

最有意思的事情就发生在这些摇摆之中。白天出现过的神经元触发序列会重新“回放”——而且快好几倍：在记忆脑区大约快二十倍，在皮层快几倍。如果人为增强这些回放与慢速摇摆的同步，记忆就会变好——这是一个因果性实验。这有什么用？一个可信的假说是：缓慢的长期记忆靠把新内容与旧内容交错重复来学习，免得新的覆盖掉旧的。工程师熟悉这个麻烦：一个在新任务上训练的神经网络，如果不重复旧的样本，就会忘掉旧任务。

\section*{修剪连接}

另一件已测量到的事：睡眠之后，神经元之间的接点平均会变小——按各自的大小成比例地缩小几个到几十个百分点。就好像把所有音量一下子统一调低了一点：比例保持不变，学到的东西还在，而明天学习所需的空间腾出来了。至于这是不是睡眠的主要任务——还是假说。

\section*{梦与清扫}

做梦的睡眠像一次试验台测试：机器全功率运转，但输入被换成了内部信号，输出被断开。它有什么用，我们不知道。至于睡眠时大脑会把废物“冲洗”掉这个想法，目前还有争议：支持和反对的实验都有。

最令人惊讶的是：睡眠可以是局部的。长时间不睡之后，一个清醒的人的皮层中，个别区域会在工作中途突然沉默，片刻“睡着”。

\fullversion{20}
